Symbols are used consistently across every chapter. Where the implementation
uses a different name for the same object, the code identifier is given in the
last column.

\vspace{1em}
\noindent
\begin{longtable}{@{}l l l@{}}
\toprule
\textbf{Symbol} & \textbf{Meaning} & \textbf{Code} \\
\midrule
\endfirsthead
\toprule
\textbf{Symbol} & \textbf{Meaning} & \textbf{Code} \\
\midrule
\endhead
\multicolumn{3}{@{}l}{\emph{Circuit}}\\[2pt]
$n$              & number of non-ground nodes                    & \code{circuit.node\_count} \\
$n_b$            & number of nonlinear branches                  & \code{circuit.branch\_count} \\
$\xvec(t)\in\mathbb{R}^{n}$ & node fluxes                        & \code{x\_t} \\
$\Cmat,\Gmat,\Kmat\in\mathbb{C}^{n\times n}$ & capacitance, conductance, linear stiffness & \code{circuit.C/G/K} \\
$\Bphi\in\mathbb{R}^{n\times n_b}$ & node--branch incidence of the nonlinear branches & \code{circuit.Bphi} \\
$\psi\in\mathbb{R}^{n_b}$ & branch fluxes, $\psi=\Bphi^{\mathsf T}\xvec$ & \code{psi\_t} \\
$\psi_{\mathrm{dc}}$ & static branch-flux offset                  & \code{dc\_branch\_flux} \\
$\Ic$            & branch critical currents                       & \code{circuit.Ic} \\
$\rphi$          & reduced flux quantum $\fluxquantum/2\pi$       & \code{circuit.phi0} \\
$Z_0$            & port reference impedance                      & \code{z0\_ohm} \\
\addlinespace
\multicolumn{3}{@{}l}{\emph{Pump harmonic balance}}\\[2pt]
$\wp$            & pump angular frequency                        & \code{grid.omega} \\
$\modeset$       & retained pump-mode set, $k\in\modeset$        & \code{grid.modes} \\
$H=|\modeset|$   & number of retained pump modes                 & \code{problem.H} \\
$X_k\in\mathbb{C}^{n}$ & pump phasor of mode $k$                 & \code{X[h]} \\
$N_t$            & AFT time samples per pump period              & \code{grid.nt} \\
$\Dmat_k$        & linear dynamic block at $k\wp$                & \code{\_linear\_blocks[h]} \\
$\Svec_k$        & source coefficient of mode $k$                & \code{source\_coeffs} \\
$\lambda$        & continuation (source-scale) parameter         & \code{source\_scale} \\
$\Rvec$          & harmonic-balance residual                     & \code{residual\_coeffs} \\
$\Nvec$          & projected nonlinear branch current            & \code{nonlinear\_current\_coeffs} \\
$\Jmat$          & Jacobian $\partial\Rvec/\partial\Xvec$        & \code{jvp\_coeffs} \\
$\Mmat$          & preconditioner, $\Mmat\approx\Jmat$           & \code{Mop} \\
$\gamma(t)$      & junction tangent (inverse inductance)         & \code{tangent.gamma\_t} \\
$\ghat_\ell$     & $\ell$-th Fourier coefficient of $\gamma$     & \code{gamma\_hat[ell]} \\
$\Khat_\ell$     & tangent-stiffness harmonic                    & \code{khat[ell]} \\
$\alpha$         & line-search damping factor                    & \code{alpha} \\
\addlinespace
\multicolumn{3}{@{}l}{\emph{Small-signal (Floquet) stage}}\\[2pt]
$\ws$            & signal angular frequency                      & \code{omega\_s} \\
$m$              & Floquet sideband index                        & \code{signal\_m} \\
$S$              & sideband truncation, $m\in[-S,S]$             & \code{sidebands} \\
$\omega_m$       & $\ws+m\wp$                                    & \code{omega\_m} \\
$\mat{A}$        & sideband conversion matrix                    & \code{assemble\_conversion\_matrix} \\
\addlinespace
\multicolumn{3}{@{}l}{\emph{Saturation (multitone) stage}}\\[2pt]
$\delta$         & pump--signal detuning $\wp-\ws$               & \code{basis.delta} \\
$(h,q)$          & tone lattice coordinate                       & \code{ToneIndex} \\
$\toneset$       & retained tone set                             & \code{basis.tones} \\
$N_p,N_\delta$   & torus grid sizes                              & \code{n\_p}, \code{n\_delta} \\
$\tau$           & affine source-path parameter                  & \code{tau} \\
$P_{1\mathrm{dB}}$ & input 1\,dB compression point               & \code{p1db} \\
\bottomrule
\end{longtable}

\paragraph{Two flux quanta.} Throughout, $\rphi=\fluxquantum/2\pi\approx
\SI{3.291e-16}{\weber}$ is the \emph{reduced} flux quantum. Every occurrence of
$\rphi$ in this thesis is the reduced one; the unreduced $\fluxquantum$ appears
only where an external flux bias is expressed as a fraction $\Phi_{\mathrm{ext}}
/\fluxquantum$. Confusing the two is a factor of $2\pi$ in the junction
nonlinearity and is one of the more common ways to produce a plausible but
wrong gain curve.

\paragraph{Phasor convention.} All harmonic amplitudes are
\emph{positive-frequency phasors} reconstructed with an explicit factor of two,
\[
  \xvec(t) = 2\,\Real\!\left[\sum_{k\in\modeset} X_k e^{+ik\wp t}\right].
\]
The factor $2$ and the sign $+i$ are not cosmetic: they propagate into the
source amplitude, the power conversion, and the interpretation of any pump
solution imported from another code. \Cref{chap:conventions} collects every
place this convention has an observable consequence.
